\section{Introduction}
\label{sec:intro}

Generative models have emerged as a powerful policy class for continuous
control.  Diffusion policies~\citep{chi2023diffusion} and flow-matching
policies~\citep{lipman2023flow} parameterize the action distribution as the
output of an iterative denoising or transport process, enabling expressive
multimodal action distributions that outperform simple Gaussian policies on
complex manipulation and locomotion tasks.

When such generative policies are fine-tuned with on-policy reinforcement
learning---specifically Proximal Policy Optimization
(PPO)~\citep{schulman2017ppo}---the standard practice is to treat the entire
generative chain as a \emph{black box}.  The holistic log-probability
$\log p_\theta(\mathbf{a}\mid\mathbf{s})$ is obtained by summing per-step
log-conditionals and a change-of-variable correction, and a single
importance-sampling ratio $r_t = p_\theta(\mathbf{a}_t\mid\mathbf{s}_t) /
p_{\theta_{\text{old}}}(\mathbf{a}_t\mid\mathbf{s}_t)$ drives the PPO
surrogate loss.

However, a stochastic flow policy with $K$ internal steps exposes $K$
\emph{tractable} conditional densities
$p_\theta(\mathbf{z}_{k-1}\mid\mathbf{z}_k,\mathbf{s},k)$, each a simple
Gaussian.  This raises a natural design question:

\begin{quote}
\emph{Should PPO operate holistically on the final action, or per-step on each
internal transition?  And if per-step, should every step receive equal credit
for the task reward, or should credit be weighted?}
\end{quote}

Per-step optimization is appealing because it decomposes a single, potentially
high-variance ratio into $K$ smaller, better-behaved ratios, each of which can
be individually clipped.  Weighted credit assignment further allows the
algorithm to focus gradient signal on the steps that matter most for downstream
performance.

\paragraph{Contributions.}
We make the following contributions:
\begin{enumerate}[leftmargin=*,nosep]
  \item We \textbf{formalize per-step PPO} for stochastic flow policies,
        deriving per-step importance-sampling ratios and clipped surrogate
        objectives (\Cref{sec:method}).
  \item We \textbf{propose weighted credit-assignment} schemes---uniform,
        learned-global, and state-dependent---that modulate how the task
        advantage is distributed across generative steps.
  \item We provide an \textbf{empirical comparison} on continuous-control
        benchmarks (HalfCheetah-v5, Hopper-v5), analyzing learning speed,
        final performance, clip fractions, and learned weight profiles.
\end{enumerate}
